\documentclass[10pt,a4paper]{amsart}
\usepackage{amsmath,amssymb,amsthm,graphicx}
\usepackage{verbatim}
\newenvironment{reference}{\comment}{\endcomment}
%\newenvironment{reference}{}{}
\newenvironment{slogan}{\comment}{\endcomment}
\newenvironment{history}{\comment}{\endcomment}

% For commutative diagrams we use Xy-pic
\usepackage[all]{xy}

% We use 2cell for 2-commutative diagrams.
\xyoption{2cell}
% Theorem environments.
%
\theoremstyle{plain}
\newtheorem{theorem}[subsection]{Theorem}
\newtheorem{proposition}[subsection]{Proposition}
\newtheorem{lemma}[subsection]{Lemma}

\theoremstyle{definition}
\newtheorem{definition}[subsection]{Definition}
\newtheorem{example}[subsection]{Example}
\newtheorem{exercise}[subsection]{Exercise}
\newtheorem{situation}[subsection]{Situation}

\theoremstyle{remark}
\newtheorem{remark}[subsection]{Remark}
\newtheorem{remarks}[subsection]{Remarks}

\numberwithin{equation}{subsection}

% Macros
%
\def\lim{\mathop{\mathrm{lim}}\nolimits}
\def\colim{\mathop{\mathrm{colim}}\nolimits}
\def\Spec{\mathop{\mathrm{Spec}}}
\def\Hom{\mathop{\mathrm{Hom}}\nolimits}
\def\Ext{\mathop{\mathrm{Ext}}\nolimits}
\def\SheafHom{\mathop{\mathcal{H}\!\mathit{om}}\nolimits}
\def\SheafExt{\mathop{\mathcal{E}\!\mathit{xt}}\nolimits}
\def\Sch{\mathit{Sch}}
\def\Mor{\mathop{\mathrm{Mor}}\nolimits}
\def\Ob{\mathop{\mathrm{Ob}}\nolimits}
\def\Sh{\mathop{\mathit{Sh}}\nolimits}
\def\NL{\mathop{N\!L}\nolimits}
\def\CH{\mathop{\mathrm{CH}}\nolimits}
\def\proetale{{pro\text{-}\acute{e}tale}}
\def\etale{{\acute{e}tale}}
\def\QCoh{\mathit{QCoh}}
\def\Ker{\mathop{\mathrm{Ker}}}
\def\Im{\mathop{\mathrm{Im}}}
\def\Coker{\mathop{\mathrm{Coker}}}
\def\Coim{\mathop{\mathrm{Coim}}}

% Boxtimes
%
\DeclareMathSymbol{\boxtimes}{\mathbin}{AMSa}{"02}

%
% Macros for moduli stacks/spaces
%
\def\QCohstack{\mathcal{QC}\!\mathit{oh}}
\def\Cohstack{\mathcal{C}\!\mathit{oh}}
\def\Spacesstack{\mathcal{S}\!\mathit{paces}}
\def\Quotfunctor{\mathrm{Quot}}
\def\Hilbfunctor{\mathrm{Hilb}}
\def\Curvesstack{\mathcal{C}\!\mathit{urves}}
\def\Polarizedstack{\mathcal{P}\!\mathit{olarized}}
\def\Complexesstack{\mathcal{C}\!\mathit{omplexes}}
% \Pic is the operator that assigns to X its picard group, usage \Pic(X)
% \Picardstack_{X/B} denotes the Picard stack of X over B
% \Picardfunctor_{X/B} denotes the Picard functor of X over B
\def\Pic{\mathop{\mathrm{Pic}}\nolimits}
\def\Picardstack{\mathcal{P}\!\mathit{ic}}
\def\Picardfunctor{\mathrm{Pic}}
\def\Deformationcategory{\mathcal{D}\!\mathit{ef}}
\usepackage{cleveref}
\crefname{originallemma}{original lemma}{original lemmas}
\Crefname{originallemma}{Original lemma}{Original lemmas}
\makeatletter
\newlabel{derived-lemma-subcategory-left-resolution}{{05T7}{source}{}{sourcepointer-1}{}}
\newlabel{derived-lemma-subcategory-left-resolution@cref}{{[originallemma][0][]05T7}{[1][0][]source}{}{}{}}
\makeatother
\crefname{originallemma}{original lemma}{original lemmas}
\Crefname{originallemma}{Original lemma}{Original lemmas}
\makeatletter
\newlabel{derived-lemma-subcategory-left-acyclics}{{05T9}{source}{}{sourcepointer-2}{}}
\newlabel{derived-lemma-subcategory-left-acyclics@cref}{{[originallemma][0][]05T9}{[1][0][]source}{}{}{}}
\makeatother
\crefname{originalproposition}{original proposition}{original propositions}
\Crefname{originalproposition}{Original proposition}{Original propositions}
\makeatletter
\newlabel{derived-proposition-enough-acyclics}{{05TA}{source}{}{sourcepointer-3}{}}
\newlabel{derived-proposition-enough-acyclics@cref}{{[originalproposition][0][]05TA}{[1][0][]source}{}{}{}}
\makeatother
\crefname{originallemma}{original lemma}{original lemmas}
\Crefname{originallemma}{Original lemma}{Original lemmas}
\makeatletter
\newlabel{derived-lemma-pre-derived-adjoint-functors}{{0DVC}{source}{}{sourcepointer-4}{}}
\newlabel{derived-lemma-pre-derived-adjoint-functors@cref}{{[originallemma][0][]0DVC}{[1][0][]source}{}{}{}}
\makeatother
\title{Group cohomology of modules with continuous action agrees with that of the underlying modules}
\author{The Stacks Project Authors}
\date{stacks-project:etale-cohomology.tex:local-snapshot-20261011; GNU FDL 1.2 or later}
\begin{document}
\maketitle
\noindent\textit{Editorial note.} All delivered mathematical text is the author's own, copied from the pinned Stacks Project chapter file etale-cohomology.tex; no statement or proof step was rewritten, reordered or removed.\par\medskip
\noindent\textit{Editorial note.} The proof is the author's complete argument for the commutativity of the diagram in all degrees.\par\medskip
\noindent\textit{Editorial note.} Every labelled reference occurring inside the counted spans is carried in assembly.external\_references as a non-load-bearing navigation pointer (same-chapter ones at their definition site, cross-chapter ones at the invocation inside the counted span), so no load-bearing mathematics is replaced by a stub. No bibliography citation occurs inside the counted spans.\par\medskip
\noindent\textit{Editorial note.} The checked chapter list for same-name collisions is the set of stacks chapters staged in this repository: chow, cohomology, divisors, etale-cohomology, flat, more-algebra, more-morphisms, morphisms, spaces-more-morphisms, varieties together with this card's chapter etale-cohomology.\par\medskip
\noindent\textit{Editorial note.} The preamble is the chapter's own native preamble excerpt already used for this source: the Xy-pic setup (native preamble lines 14-18) and the native theorem declarations with the mathematical macros (native preamble lines 154-217).\par\medskip
\noindent\textit{Allowed standard prerequisites (editorial).}\begin{itemize}
\item Category of R-G-modules and the forgetful functor to G-modules. Same chapter, the category of modules over a ring with a continuous action of a topological group and the forgetful functor forgetting the R-structure. Supplies the two vertical functors in the diagram of the statement.
\item Continuous cohomology H\textasciicircum{}i(G, -) and the comparison between the two categories. Same chapter, the definition of the continuous cohomology functors and the companion result comparing the cohomology of a module with that of its underlying module. Supplies the two horizontal functors and the comparison established in the proof.
\end{itemize}\medskip
\ifdefined\hypertarget\hypertarget{sourcepointer-1}{}\fi
\noindent\textit{Source reference.} Original lemma 05T7: sources/stacks-etale-cohomology/pinned-original-source.tex, source lines 8884--8884. Navigation pointer invoked at proof line 8884; the counted span contains the pointer only and the cited result is not reproduced inside it, so the pointer is registered as a non-load-bearing navigation stub in the established form for Stacks material. No bibliography citation is involved.\par\medskip
\ifdefined\hypertarget\hypertarget{sourcepointer-2}{}\fi
\noindent\textit{Source reference.} Original lemma 05T9: sources/stacks-etale-cohomology/pinned-original-source.tex, source lines 8888--8888. Navigation pointer invoked at proof line 8888; the counted span contains the pointer only and the cited result is not reproduced inside it, so the pointer is registered as a non-load-bearing navigation stub in the established form for Stacks material. No bibliography citation is involved.\par\medskip
\ifdefined\hypertarget\hypertarget{sourcepointer-3}{}\fi
\noindent\textit{Source reference.} Original proposition 05TA: sources/stacks-etale-cohomology/pinned-original-source.tex, source lines 8889--8889. Navigation pointer invoked at proof line 8889; the counted span contains the pointer only and the cited result is not reproduced inside it, so the pointer is registered as a non-load-bearing navigation stub in the established form for Stacks material. No bibliography citation is involved.\par\medskip
\ifdefined\hypertarget\hypertarget{sourcepointer-4}{}\fi
\noindent\textit{Source reference.} Original lemma 0DVC: sources/stacks-etale-cohomology/pinned-original-source.tex, source lines 8891--8891. Navigation pointer invoked at proof line 8891; the counted span contains the pointer only and the cited result is not reproduced inside it, so the pointer is registered as a non-load-bearing navigation stub in the established form for Stacks material. No bibliography citation is involved.\par\medskip
\begin{lemma}
\label{lemma-modules-abelian}
Let $G$ be a topological group. Let $R$ be a ring.
For every $i \geq 0$ the diagram
$$
\xymatrix{
\text{Mod}_{R, G} \ar[rr]_{H^i(G, -)} \ar[d] & &
\text{Mod}_R \ar[d] \\
\text{Mod}_G \ar[rr]^{H^i(G, -)} & &
\textit{Ab}
}
$$
whose vertical arrows are the forgetful functors is commutative.
\end{lemma}

\begin{proof}
Let us denote the forgetful functor $F : \text{Mod}_{R, G} \to \text{Mod}_G$.
Then $F$ has a left adjoint $H : \text{Mod}_G \to \text{Mod}_{R, G}$
given by $H(M) = M \otimes_\mathbf{Z} R$. Observe that every object of
$\text{Mod}_G$ is a quotient of a direct sum of modules of the form
$\mathbf{Z}[G/U]$ where $U \subset G$ is an open subgroup.
Here $\mathbf{Z}[G/U]$ denotes the $G$-modules of
finite $\mathbf{Z}$-linear combinations
of right $U$ congruence classes in $G$ endowed with left $G$-action.
Thus every bounded above complex in $\text{Mod}_G$ is quasi-isomorphic
to a bounded above complex in $\text{Mod}_G$ whose underlying
terms are flat $\mathbf{Z}$-modules
(Derived Categories, Lemma \ref{derived-lemma-subcategory-left-resolution}).
Thus it is clear that $LH$ exists on $D^-(\text{Mod}_G)$ and is computed by
evaluating $H$ on any complex whose terms are flat $\mathbf{Z}$-modules;
this follows from
Derived Categories, Lemma \ref{derived-lemma-subcategory-left-acyclics} and
Proposition \ref{derived-proposition-enough-acyclics}.
We conclude from Derived Categories, Lemma
\ref{derived-lemma-pre-derived-adjoint-functors}
that
$$
\text{Ext}^i(\mathbf{Z}, F(M)) = \text{Ext}^i(R, M)
$$
for $M$ in $\text{Mod}_{R, G}$.
Observe that $H^0(G, -) = \Hom(\mathbf{Z}, -)$ on
$\text{Mod}_G$ where $\mathbf{Z}$ denotes the $G$-module
with trivial action. Hence
$H^i(G, -) = \text{Ext}^i(\mathbf{Z}, -)$ on $\text{Mod}_G$.
Similarly we have $H^i(G, -) = \text{Ext}^i(R, -)$ on
$\text{Mod}_{R, G}$. Combining everything we see that the lemma is true.
\end{proof}

\end{document}
